\documentclass[11pt]{amsart}

\usepackage[T1]{fontenc}
\usepackage{lmodern}
\usepackage{microtype}
\usepackage{mathtools,amssymb,amsthm}
\usepackage[hidelinks]{hyperref}

\hypersetup{
  pdftitle={A mod 32 counterexample to the congruence form of Propp's
            Conjecture 2 for Aztec diamond tilings by dominoes and square
            tetrominoes}
}

\newtheorem{theorem}{Theorem}

\DeclareMathOperator{\vtwo}{v_2}

\emergencystretch=3em

\title[A mod $32$ counterexample to Propp's Conjecture 2, congruence form]
{A mod $32$ counterexample to the congruence form of Propp's Conjecture 2
for Aztec diamond tilings by dominoes and square tetrominoes}

\date{}

\subjclass[2020]{05B45, 05A15, 11A07}
\keywords{Aztec diamond, polyomino tiling, square tetromino, 2-adic valuation,
congruence, transfer matrix, counterexample, OEIS A356512}

\begin{document}

\begin{abstract}
Let $M(n)$ be the number of tilings of the Aztec diamond of order $n$ by
dominoes and square tetrominoes (OEIS A356512). Conjecture~2 of Propp
[\emph{Integers} \textbf{23} (2023), Art.~A30] asserts that for all $k\ge1$,
$n+n'\equiv-3 \pmod{2^k}$ implies $M(n)+M(n')\equiv0 \pmod{2^k}$. We refute this congruence form
at the single instance $(n,n',k)=(14,15,5)$: here $14+15=29\equiv-3
\pmod{32}$, while
\[
  M(14)\equiv 79,\qquad M(15)\equiv 225 \pmod{256},
\]
so $M(14)+M(15)\equiv48\pmod{256}$, whose $2$-adic valuation is $4$, not
$\ge5$. Propp's Conjecture~1, on the $2$-adic continuity of $M$, is not
refuted, and neither is the antisymmetry restatement of Conjecture~2, which
presupposes it.
\end{abstract}

\maketitle

\section{The conjecture}\label{sec:conj}

For $n\ge0$ let $AD_n$ be the Aztec diamond region of order $n$: the union of
the closed unit squares of the integer lattice that lie entirely inside
$|x|+|y|\le n+1$. Equivalently $AD_n$ has $2n$ rows, the row $r$
($0\le r\le 2n-1$) consisting of $2\min(r+1,\,2n-r)$ cells centred on the
vertical axis, so that
\begin{equation}\label{eq:cells}
  |AD_n| \;=\; 2n(n+1).
\end{equation}
Let
\[
  M(n) \;=\; \#\bigl\{\text{tilings of } AD_n \text{ by } 1\times2,\ 2\times1
  \text{ and } 2\times2 \text{ tiles}\bigr\},
\]
that is, by dominoes in either orientation together with the square
tetromino. Then $M(0)=1$, $M(1)=3$, and $M$ is the sequence
OEIS~\textbf{A356512}. Propp \cite{Propp} states two conjectures about the
$2$-adic behaviour of $M$; the second, quoted here from p.~6 of the published
version, is the object of this note.

\begin{quote}
\textbf{Conjecture 2} (Propp \cite{Propp}). \emph{For all $k\ge1$, if
$n+n'\equiv-3 \pmod{2^k}$ then $M(n)+M(n')\equiv0 \pmod{2^k}$.}
\end{quote}

Both $n,n'$ range over $\mathbb{Z}_{\ge0}$ and $k$ over $\mathbb{Z}_{\ge1}$;
there is no existential quantifier, so a single triple $(n,n',k)$ satisfying
the hypothesis and violating the conclusion refutes the statement. Propp
prints $M(n)$ for $0\le n\le12$ and remarks that he has no efficient method of
computing further terms; since $n,n'\le12$ forces $n+n'\le24<29$, that table
cannot test $k=5$ at any point.

\section{The counterexample}\label{sec:refute}

\begin{theorem}\label{thm:main}
The congruence form of Conjecture~$2$ of \cite{Propp}, as quoted in
\S\ref{sec:conj}, is false. Explicitly, take $n=14$, $n'=15$ and
$k=5$. Then $n+n'\equiv-3 \pmod{2^k}$, while
$\vtwo\bigl(M(14)+M(15)\bigr)=4$ and so $M(n)+M(n')\not\equiv0 \pmod{2^{k}}$.
\end{theorem}

\begin{proof}
The two residues, computed as described in \S\ref{sec:verify}, are
\begin{equation}\label{eq:M1415}
  M(14)\equiv 79 \pmod{256},
  \qquad
  M(15)\equiv 225 \pmod{256}.
\end{equation}
The hypothesis holds: $14+15=29$ and $29+3=32=2^5$, so $14+15\equiv-3
\pmod{32}$; equivalently $(-3-14) \bmod 32 = 15 = n'$. Both indices are
nonnegative and $k=5\ge1$, so the triple lies in the range quantified over.

The conclusion fails. By \eqref{eq:M1415}, $M(14)+M(15)\equiv 79+225=304\equiv
48\pmod{256}$, and $48=2^4\cdot3$. A residue modulo $256=2^8$ determines the
$2$-adic valuation of an integer whenever that valuation is less than $8$, so
$\vtwo(M(14)+M(15))=\vtwo(48)=4$. In particular
$M(14)+M(15)\equiv16\not\equiv0\pmod{32}$. The valuation is short of the
required $5$ by exactly one power of two.
\end{proof}

\section{The computation and its controls}\label{sec:verify}

The residues \eqref{eq:M1415} were obtained by a dense broken-profile transfer
computation over the $2n\times2n$ bounding box of $AD_n$, one cell at a time,
with $2n$ frontier bits and one extra bit for the cell that a $2\times2$ tile
placed at the previous column occupies on the next row; cells outside $AD_n$
are treated as permanently blocked, and the programs assert that the number of
admissible cells equals $2n(n+1)$ as in \eqref{eq:cells}. At $n=14$ and $n=15$
this carries $2^{29}$ and $2^{31}$ states. Two implementations written
independently of each other, in C, agreed on both residues, and each residue
was produced twice, on different machines. A separate corroboration of $M(14)$
is arithmetic: an exact accumulator modulo $2^{128}$ returned
$M(14)\equiv 228062845944308165005028192227771743567 \pmod{2^{128}}$, and
since $2^8$ divides $10^8$ the last eight decimal digits of that number
determine the residue modulo $256$, giving $71743567 = 280248\cdot256+79$, in
agreement with \eqref{eq:M1415}.

A third implementation accompanies this note as \texttt{verify.py}: a pure
Python program, standard library only, exact integer arithmetic, sharing no
code with either C engine. It recomputes $AD_n$ from the inequality
$|x|+|y|\le n+1$ and checks \eqref{eq:cells}; reproduces $M(0),\dots,M(4)$ by
plain backtracking with no transfer computation at all; reproduces all thirteen
published terms of A356512 exactly; reproduces the domino-only specialisation
$2^{n(n+1)/2}$ of Elkies, Kuperberg, Larsen and Propp \cite{EKLP}
(OEIS A006125) for $n\le12$ as a control on the region and the domino rules;
then recomputes $M(14)$ and $M(15) \bmod 256$ from scratch, obtaining $79$ and
$225$, and re-derives Theorem~\ref{thm:main} from its own residues rather than
from the printed ones. Its recorded run, \texttt{verify.output.txt}, reports
\texttt{VERDICT: ALL 63 CHECKS PASS}. The deep part is expensive for an
interpreted language: in that run the two residues took $777$ and $1835$
seconds on one core and the process peaked at $9.81$ GB.

Two remarks of that transcript overreach --- its check \texttt{E8} and its
closing scope note --- and nothing here relies on either.

\section{Scope}\label{sec:scope}

Only the congruence form of Conjecture~2, exactly as printed and quoted in
\S\ref{sec:conj}, is refuted. Propp restates it as: \emph{``if one extends
$M:\mathbb{N}\to\mathbb{N}$ to the $2$-adic function
$\widehat M:\mathbb{Z}_2\to\mathbb{Z}_2$, one has
$\widehat M(-3-n)=-\widehat M(n)$''}; that form presupposes the $2$-adic
extension, i.e.\ his Conjecture~1, which is neither used nor refuted here.
$M(14)$ and $M(15)$ are established modulo $256$ only; the exact integers are
not computed anywhere here, and Theorem~\ref{thm:main} does not need them.
Nothing is claimed for an index $\ge16$: every other pair on the diagonal
$n+n'=29$ has an index $\ge16$, where $M$ is unknown to us, so those instances
are settled in neither direction, no instance with $k\ge6$ was tested, and
$(14,15,5)$ is not claimed to be the least failure. Finally, every control
listed in \S\ref{sec:verify} lies at $n\le12$, the range where published data
exists; at $n=14$ and $n=15$ there is nothing external to compare against, and
what stands in place of a control there is the agreement of the three
independently written programs.

\begin{thebibliography}{9}

\bibitem{EKLP}
N.~Elkies, G.~Kuperberg, M.~Larsen, and J.~Propp,
\emph{Alternating-sign matrices and domino tilings I},
J. Algebraic Combin. \textbf{1} (1992), no.~2, 111--132.
\href{https://doi.org/10.1023/A:1022420103267}{doi:10.1023/A:1022420103267}.
The specialisation used as a control here, $2^{n(n+1)/2}$ domino tilings of
$AD_n$, is OEIS A006125.

\bibitem{OEIS}
OEIS Foundation Inc.,
\emph{The On-Line Encyclopedia of Integer Sequences},
sequence \href{https://oeis.org/A356512}{A356512} ($M$, contributed by
J.~Propp, August 2022; thirteen terms, $n=0,\dots,12$) and
\href{https://oeis.org/A006125}{A006125}.

\bibitem{Propp}
J.~Propp,
\emph{Some $2$-adic conjectures concerning polyomino tilings of Aztec
diamonds},
Integers \textbf{23} (2023), Art.~A30;
\href{https://arxiv.org/abs/2204.00158}{arXiv:2204.00158}.
Conjecture~2 is quoted from the published version, p.~6.

\end{thebibliography}

\end{document}
